% GENERATED FILE. Edit canonical lessons, then run npm run generate:books.
% canonical-source-hash: fnv1a64:87a189f5a4ac6a64
% canonical-lessons: SA-C324-pura, SA-C324-itivrttam, SA-C324-puranam, SA-C324-anusmarati, SA-C324-kirtayati

\chapter{Tales of History}
\label{ch:sa-tales-of-history}
\hlchaptermodality{324}

\begin{chapteropening}
\textbf{By the end of this chapter:} \emph{I can tell what happened long ago and how it is remembered.}

Complete the last lesson of chapter 324: I can tell what happened long ago and how it is remembered.
\end{chapteropening}

\section[\texorpdfstring{\sk{पुरा} (purā)}{पुरा (purā)}]{\sk{पुरा} (purā) — long ago; formerly}

\label{lesson:SA-C324-pura}



\begin{quote}
Before the new one: say the Sanskrit for digs up, then the Sanskrit for brings to light.
\end{quote}

\subsection*{You'll want to know: \sk{पुरा}}
\textbf{\sk{पुरा}} — \emph{purā} — \textquotedblleft{}long ago; formerly\textquotedblright{}.

\begin{quote}
\sk{पुरा} \sk{अत्र} \sk{एकः} \sk{राजा} \sk{वसति} \sk{स्म।}
\end{quote}

\begin{quote}
\emph{purā atra ekaḥ rājā vasati sma.}
\end{quote}

\begin{quote}
Long ago a king lived here.
\end{quote}

\begin{cousinweb}
\textbf{\sk{पुरा}} (\emph{purā}) MW gives \emph{purā} as before, formerly, of old, and compares \emph{puras} and \emph{pūrva}, but gives no root, so its origin is left uncertain here (MW).
\end{cousinweb}

\begin{grammarlens}[title={What you've built}]
One more word for telling what happened.
\end{grammarlens}

\subsection*{Your turn}
\begin{itemize}
\raggedright
  \item \emph{Say it:} \emph{purā}
  \item \emph{Say it:} \emph{purā}, once more
  \item \emph{Say it:} \emph{purā atra ekaḥ rājā vasati sma.} — Long ago a king lived here.
  \item \emph{Recall:} say the Sanskrit for digs up, then the Sanskrit for brings to light, then say \emph{purā} again
\end{itemize}

\begin{tcolorbox}[breakable,colback=teal!4,colframe=teal!35!black,title={Before you move on}]
What does \sk{पुरा} mean? (\textquotedblleft{}Long ago; formerly\textquotedblright{}.) Say it once more.
\end{tcolorbox}
\section[\texorpdfstring{\sk{इतिवृत्तम्} (itivṛttam)}{इतिवृत्तम् (itivṛttam)}]{\sk{इतिवृत्तम्} (itivṛttam) — a chronicle; an account of events}

\label{lesson:SA-C324-itivrttam}



\begin{quote}
Before the new one: say the Sanskrit for brings to light, then the Sanskrit for long ago.
\end{quote}

\subsection*{You'll want to know: \sk{इतिवृत्तम्}}
\textbf{\sk{इतिवृत्तम्}} — \emph{itivṛttam} — \textquotedblleft{}a chronicle; an account of events\textquotedblright{}.

\begin{quote}
\sk{एतत्} \sk{पुस्तकं} \sk{राज्यस्य} \sk{इतिवृत्तं} \sk{वर्णयति।}
\end{quote}

\begin{quote}
\emph{etat pustakaṃ rājyasya itivṛttaṃ varṇayati.}
\end{quote}

\begin{quote}
This book tells the chronicle of the kingdom.
\end{quote}

A neuter noun.

\begin{cousinweb}
\textbf{\sk{इतिवृत्तम्}} (\emph{itivṛttam}) is built from \emph{iti}, \textquotedblleft{}thus\textquotedblright{}, and \emph{vṛtta}, \textquotedblleft{}happened\textquotedblright{}: what happened thus; MW gives \emph{itivṛtta} as an occurrence, an event, and a chronicle is its modern use (MW).
\end{cousinweb}

\begin{grammarlens}[title={What you've built}]
One more word for telling what happened.
\end{grammarlens}

\subsection*{Your turn}
\begin{itemize}
\raggedright
  \item \emph{Say it:} \emph{itivṛttam}
  \item \emph{Say it:} \emph{itivṛttam}, once more
  \item \emph{Say it:} \emph{etat pustakaṃ rājyasya itivṛttaṃ varṇayati.} — This book tells the chronicle of the kingdom.
  \item \emph{Recall:} say the Sanskrit for brings to light, then the Sanskrit for long ago, then say \emph{itivṛttam} again
\end{itemize}

\begin{tcolorbox}[breakable,colback=teal!4,colframe=teal!35!black,title={Before you move on}]
What does \sk{इतिवृत्तम्} mean? (\textquotedblleft{}A chronicle; an account of events\textquotedblright{}.) Say it once more.
\end{tcolorbox}
\section[\texorpdfstring{\sk{पुराणम्} (purāṇam)}{पुराणम् (purāṇam)}]{\sk{पुराणम्} (purāṇam) — an ancient tale; a Purana}

\label{lesson:SA-C324-puranam}



\begin{quote}
Before the new one: say the Sanskrit for long ago, then the Sanskrit for a chronicle.
\end{quote}

\subsection*{You'll want to know: \sk{पुराणम्}}
\textbf{\sk{पुराणम्}} — \emph{purāṇam} — \textquotedblleft{}an ancient tale; a Purana\textquotedblright{}.

\begin{quote}
\sk{पितामहः} \sk{पुराणं} \sk{कथयति।}
\end{quote}

\begin{quote}
\emph{pitāmahaḥ purāṇaṃ kathayati.}
\end{quote}

\begin{quote}
Grandfather tells an ancient tale.
\end{quote}

A neuter noun.

\begin{cousinweb}
\textbf{\sk{पुराणम्}} (\emph{purāṇam}) is usually traced to \emph{purā}, \textquotedblleft{}formerly\textquotedblright{}; MW gives \emph{purāṇa} as belonging to olden times and, as a noun, an ancient tale or legend, old traditional history (MW).
\end{cousinweb}

\begin{grammarlens}[title={What you've built}]
One more word for telling what happened.
\end{grammarlens}

\subsection*{Your turn}
\begin{itemize}
\raggedright
  \item \emph{Say it:} \emph{purāṇam}
  \item \emph{Say it:} \emph{purāṇam}, once more
  \item \emph{Say it:} \emph{pitāmahaḥ purāṇaṃ kathayati.} — Grandfather tells an ancient tale.
  \item \emph{Recall:} say the Sanskrit for long ago, then the Sanskrit for a chronicle, then say \emph{purāṇam} again
\end{itemize}

\begin{tcolorbox}[breakable,colback=teal!4,colframe=teal!35!black,title={Before you move on}]
What does \sk{पुराणम्} mean? (\textquotedblleft{}An ancient tale; a Purana\textquotedblright{}.) Say it once more.
\end{tcolorbox}
\section[\texorpdfstring{\sk{अनुस्मरति} (anusmarati)}{अनुस्मरति (anusmarati)}]{\sk{अनुस्मरति} (anusmarati) — calls to mind; recollects}

\label{lesson:SA-C324-anusmarati}



\begin{quote}
Before the new one: say the Sanskrit for a chronicle, then the Sanskrit for an ancient tale.
\end{quote}

\subsection*{You'll want to know: \sk{अनुस्मरति}}
\textbf{\sk{अनुस्मरति}} — \emph{anusmarati} — \textquotedblleft{}calls to mind; recollects\textquotedblright{}.

\begin{quote}
\sk{जनाः} \sk{वीरस्य} \sk{शौर्यम्} \sk{अनुस्मरन्ति।}
\end{quote}

\begin{quote}
\emph{janāḥ vīrasya śauryam anusmaranti.}
\end{quote}

\begin{quote}
People call to mind the hero's valour.
\end{quote}

Past (he …): \sk{सः} \sk{पितामहम्} \sk{अनुस्मृतवान्}, \emph{saḥ pitāmaham anusmṛtavān}.

\begin{cousinweb}
\textbf{\sk{अनुस्मरति}} (\emph{anusmarati}) is built from \emph{anu}, \textquotedblleft{}after\textquotedblright{}, and the root √smṛ, \textquotedblleft{}to remember\textquotedblright{}, as in \emph{smarati}; MW gives \emph{anu-smṛ} as to remember, recollect (MW).
\end{cousinweb}

\begin{grammarlens}[title={What you've built}]
One more word for telling what happened.
\end{grammarlens}

\subsection*{Your turn}
\begin{itemize}
\raggedright
  \item \emph{Say it:} \emph{anusmarati}
  \item \emph{Say it:} \emph{anusmarati}, once more
  \item \emph{Say it:} \emph{janāḥ vīrasya śauryam anusmaranti.} — People call to mind the hero's valour.
  \item \emph{Recall:} say the Sanskrit for a chronicle, then the Sanskrit for an ancient tale, then say \emph{anusmarati} again
\end{itemize}

\begin{tcolorbox}[breakable,colback=teal!4,colframe=teal!35!black,title={Before you move on}]
What does \sk{अनुस्मरति} mean? (\textquotedblleft{}Calls to mind; recollects\textquotedblright{}.) Say it once more.
\end{tcolorbox}
\section[\texorpdfstring{\sk{कीर्तयति} (kīrtayati)}{कीर्तयति (kīrtayati)}]{\sk{कीर्तयति} (kīrtayati) — tells of; celebrates}

\label{lesson:SA-C324-kirtayati}



\begin{quote}
Before the new one: say the Sanskrit for an ancient tale, then the Sanskrit for calls to mind.
\end{quote}

\subsection*{You'll want to know: \sk{कीर्तयति}}
\textbf{\sk{कीर्तयति}} — \emph{kīrtayati} — \textquotedblleft{}tells of; celebrates\textquotedblright{}.

\begin{quote}
\sk{ग्रन्थः} \sk{तस्य} \sk{युद्धस्य} \sk{शौर्यं} \sk{कीर्तयति।}
\end{quote}

\begin{quote}
\emph{granthaḥ tasya yuddhasya śauryaṃ kīrtayati.}
\end{quote}

\begin{quote}
The book tells of the valour shown in that war.
\end{quote}

Past (he …): \sk{सः} \sk{वीरस्य} \sk{शौर्यं} \sk{कीर्तितवान्}, \emph{saḥ vīrasya śauryaṃ kīrtitavān}.

\begin{cousinweb}
\textbf{\sk{कीर्तयति}} (\emph{kīrtayati}) comes from the root √kīrt, \textquotedblleft{}to mention, tell, name, celebrate\textquotedblright{}, which belongs with \emph{kīrti}, \textquotedblleft{}fame\textquotedblright{} (MW).
\end{cousinweb}

\begin{grammarlens}[title={What you've built}]
One more word for telling what happened.
\end{grammarlens}

\subsection*{Your turn}
\begin{itemize}
\raggedright
  \item \emph{Say it:} \emph{kīrtayati}
  \item \emph{Say it:} \emph{kīrtayati}, once more
  \item \emph{Say it:} \emph{granthaḥ tasya yuddhasya śauryaṃ kīrtayati.} — The book tells of the valour shown in that war.
  \item \emph{Recall:} say the Sanskrit for an ancient tale, then the Sanskrit for calls to mind, then say \emph{kīrtayati} again
\end{itemize}

\begin{tcolorbox}[breakable,colback=teal!4,colframe=teal!35!black,title={Before you move on}]
What does \sk{कीर्तयति} mean? (\textquotedblleft{}Tells of; celebrates\textquotedblright{}.) Say it once more.
\end{tcolorbox}
